\begin{table*}[t]
\centering
\small
\setlength{\tabcolsep}{4pt}
\caption{OhioT1DM official test split, PH = 30 min (12 subjects, $n = 28335$ windows). Coverage and width are split-conformal at $\alpha = 0.1$, calibrated on the disjoint calibration split; hypo coverage is the Mondrian $<70$ mg/dL stratum. Persistence, Ridge-AR and ARIMA have no classification head, so their PR-AUC, Se@FPR and FAPD entries are \emph{lower bounds} (Appendix~\ref{app:benchmarks}).}
\label{tab:sota_ph30}
\resizebox{\ifdim\width>\linewidth\linewidth\else\width\fi}{!}{%
\begin{tabular}{lrrrrrrrrrrrrrrr}
\toprule
Model & RMSE & MAE & MARD (\%) & Lag (min) & A+B (\%) & D+E ($n$) & Parkes A+B (\%) & PR-AUC & Se@5\% & Se@10\% & FAPD@5\% & Onset lead (min) & Cov. & Hypo cov. & Width \\
\midrule
Persistence & 23.29 & 16.79 & 11.25 & 30.0 & 99.02 & 272 & 99.66 & 0.832 & 0.933 & 0.970 & 1.16 & 23.7 & 0.913 & \textbf{0.934} & 80.0 \\
Ridge-AR & 20.05 & 14.29 & 9.76 & 25.0 & 98.74 & 351 & 99.76 & 0.829 & 0.964 & 0.990 & 2.99 & 26.4 & 0.916 & 0.933 & 67.7 \\
ARIMA(4,1,0) & 21.09 & 14.89 & 9.94 & 25.0 & 99.31 & 187 & 99.79 & 0.856 & 0.957 & 0.985 & 2.05 & 25.8 & 0.914 & 0.933 & 72.5 \\
LightGBM & \textbf{18.92} & 13.07 & 8.74 & \textbf{20.0} & 99.04 & 270 & 99.81 & 0.841 & \textbf{0.968} & 0.993 & 2.90 & \textbf{27.1} & 0.920 & 0.922 & 66.6 \\
Dilated LSTM (CGM) & 19.52 & 13.39 & 8.80 & 20.0 & \textbf{99.38} & \textbf{175} & \textbf{99.86} & 0.858 & 0.962 & \textbf{0.994} & 2.88 & 26.2 & 0.918 & 0.933 & 68.5 \\
CausalCGM-MT (CGM) & 19.16 & 13.33 & 8.89 & 20.0 & 99.18 & 231 & 99.79 & 0.868 & 0.954 & 0.990 & 2.57 & 25.6 & 0.922 & 0.926 & 68.7 \\
CausalCGM-MT (fusion) & 18.97 & \textbf{13.01} & \textbf{8.52} & 20.0 & 99.34 & 184 & 99.84 & \textbf{0.870} & 0.957 & 0.991 & 2.80 & 26.5 & 0.926 & 0.887 & 68.8 \\
CausalCGM-MT + LoRA ($r = 4$) & 21.81 & 15.82 & 10.52 & 25.0 & 98.86 & 323 & 99.67 & 0.649 & 0.783 & 0.898 & 1.80 & 21.9 & 0.921 & 0.907 & 78.3 \\
\bottomrule
\end{tabular}}
\end{table*}
